% 本文件由 paper/generate-content.js 从游戏数据自动生成，请勿手工编辑。

\subsection{赵敏线：朔漠玫瑰}

\subsubsection{第一拍·绿柳初逢}\label{question:zhaomin:1}
\textbf{情节位置：}\hyperref[plot:zhaomin:1]{绿柳初逢完整剧情}。

\textbf{题干：}赵敏所指的足底要穴——位于足底前1/3凹陷处，为足少阴肾经井穴，此穴名为：

\textbf{选项与答案}
\begin{enumerate}[label=\Alph*.,leftmargin=*]
  \item 涌泉穴\textbf{（正确答案）}
  \item 太冲穴
  \item 三阴交
  \item 足三里
\end{enumerate}
\textbf{判定依据、推导与科普边界}
\begin{quote}涌泉穴（KI1）位于足底前1/3凹陷处，是足少阴肾经的井穴。《黄帝内经》记载其为肾经所出，故名涌泉。传统针灸学常用其引火归元、醒神安神、滋肾降逆；日常按揉可作为放松与保健辅助。需注意：足底按摩、穴位按揉不能替代对月经异常、失眠、腰膝酸软等问题的正规诊疗，相关症状应结合病史、体征和现代检查辨证处理。\end{quote}
\noindent\textit{资料依据见本页脚注。}
\footnote{佚名. 黄帝内经·灵枢·本输[M/OL]. 中国哲学书电子化计划, [2026-06-18]. \href{https://ctext.org/huangdi-neijing/ben-shu/zhs}{在线链接}.}
\footnote{World Health Organization Regional Office for the Western Pacific. WHO standard acupuncture point locations in the Western Pacific region[S/OL]. Manila: WHO Regional Office for the Western Pacific, 2008[2026-06-18]. \href{https://iris.who.int/items/f188654a-d8a7-4519-9979-8e2de713c060}{在线链接}.}
\footnote{国家市场监督管理总局, 国家标准化管理委员会. GB/T 12346-2021 经穴名称与定位[S/OL]. 2021-11-26[2026-06-18]. \href{https://openstd.samr.gov.cn/bzgk/std/newGbInfo?hcno=397548AE7248D3D87DD15E0AB8107185}{在线链接}.}

\subsubsection{第二拍·黑玉断续}\label{question:zhaomin:2}
\textbf{情节位置：}\hyperref[plot:zhaomin:2]{黑玉断续完整剧情}。

\textbf{题干：}中医认为肾主骨，女性绝经后骨质疏松风险显著增加。以下关于女性骨骼健康的说法，哪一项是错误的？

\textbf{选项与答案}
\begin{enumerate}[label=\Alph*.,leftmargin=*]
  \item 绝经后雌激素下降，骨吸收大于骨形成，导致骨质疏松
  \item 中医补肾壮骨常用杜仲、续断、骨碎补等药材
  \item 骨质疏松是正常衰老现象，无需特别预防\textbf{（正确答案）}
  \item 适量负重运动、补充钙质与维生素D有助于防治骨质疏松
\end{enumerate}
\textbf{判定依据、推导与科普边界}
\begin{quote}骨质疏松并非无需预防的正常衰老现象。中国2018年流行病学调查显示，50岁以上女性骨质疏松患病率为32.1\%。绝经后雌激素下降会加快骨吸收，但预防与治疗必须个体化：优先保证膳食钙、维生素D和蛋白质摄入，进行规律的负重、抗阻与平衡训练，避免吸烟和过量饮酒；是否需要补充剂、骨密度（DXA）检查或抗骨质疏松药物，应结合年龄、骨折史和风险评估由医师决定。中医药可在正规诊疗中辨证作为辅助，但不能替代骨折风险评估和规范治疗。\end{quote}
\noindent\textit{资料依据见本页脚注。}
\footnote{中华医学会骨质疏松和骨矿盐疾病分会. 原发性骨质疏松症诊疗指南(2022)[J]. 中国全科医学, 2023, 26(14): 1671-1691. DOI：\nolinkurl{10.12114/j.issn.1007-9572.2023.0121}.}
\footnote{国家卫生健康委员会. 一图读懂: 中国骨质疏松症流行病学调查结果[R/OL]. 2018-10-20[2026-06-18]. \href{https://www.nhc.gov.cn/wjw/zcjd/201810/1aee4c12bec7415cbebfe7eacbc09552.shtml}{在线链接}.}
\footnote{Lee D O, Hong Y H, Cho M K, et al. The 2024 guidelines for osteoporosis - Korean Society of Menopause: Part I[J]. Journal of Menopausal Medicine, 2024, 30(1): 1-23. DOI：\nolinkurl{10.6118/jmm.24000}.}
\footnote{The North American Menopause Society. The 2022 hormone therapy position statement of The North American Menopause Society[J]. Menopause, 2022, 29(7): 767-794. DOI：\nolinkurl{10.1097/GME.0000000000002028}.}

\subsubsection{第三拍·万安寺火}\label{question:zhaomin:3}
\textbf{情节位置：}\hyperref[plot:zhaomin:3]{万安寺火完整剧情}。

\textbf{题干：}赵敏在火场中提醒烧伤急救步骤。现代医学关于烧伤现场急救，以下哪项处理不正确？

\textbf{选项与答案}
\begin{enumerate}[label=\Alph*.,leftmargin=*]
  \item 尽快用清洁流动水冷却创面（通常5-20分钟，遵当地急救指南）
  \item 在创面上涂抹牙膏、酱油等以缓解疼痛\textbf{（正确答案）}
  \item 小心脱去创面附近的衣物和饰品
  \item 用干净敷料覆盖创面后尽快就医
\end{enumerate}
\textbf{判定依据、推导与科普边界}
\begin{quote}在烧伤创面上涂抹牙膏、酱油、面粉等是常见的民间误区！这些物质不仅不能治疗烧伤，还会：（1）增加感染风险；（2）影响医生对创面深度的判断；（3）加重组织损伤。正确急救步骤：（1）冲——尽快用清洁、凉而不冰的流动水冷却创面，常用建议为5-20分钟，具体遵当地急救指南；（2）脱——小心去除不粘连创面的衣物和饰品；（3）盖——用干净敷料覆盖；（4）送——大面积、深部、面部/会阴/手足烧伤或儿童烧伤应及时就医。中医外治烧烫伤需在专业医师指导下使用，不宜自行敷药遮掩创面。\end{quote}
\noindent\textit{资料依据见本页脚注。}
\footnote{Hewett Brumberg E K, Douma M J, Alibertis K, et al. 2024 American Heart Association and American Red Cross guidelines for first aid[J]. Circulation, 2024, 150(24): e519-e579. DOI：\nolinkurl{10.1161/CIR.0000000000001281}.}
\footnote{World Health Organization. Burns[EB/OL]. [2026-06-18]. \href{https://www.who.int/news-room/fact-sheets/detail/burns}{在线链接}.}
\footnote{American Red Cross. Burns: types, symptoms, and how to help[EB/OL]. [2026-06-18]. \href{https://www.redcross.org/take-a-class/resources/learn-first-aid/burns}{在线链接}.}
\footnote{Mayo Clinic. Burns: first aid[EB/OL]. [2026-06-18]. \href{https://www.mayoclinic.org/first-aid/first-aid-burns/basics/art-20056649}{在线链接}.}

\subsubsection{第五拍·弥勒佛庙}\label{question:zhaomin:5}
\textbf{情节位置：}\hyperref[plot:zhaomin:5]{弥勒佛庙完整剧情}。

\textbf{题干：}张无忌提到的外伤止血与赵敏提及的月经过多（中医称崩漏），均属出血证。以下关于止血与崩漏的说法，哪一项是正确的？

\textbf{选项与答案}
\begin{enumerate}[label=\Alph*.,leftmargin=*]
  \item 崩漏只需休息即可自愈，无需治疗
  \item 中医治疗崩漏遵循塞流、澄源、复旧三法\textbf{（正确答案）}
  \item 外伤出血时，应立即在伤口上撒面粉止血
  \item 月经过多只需多喝红糖水即可
\end{enumerate}
\textbf{判定依据、推导与科普边界}
\begin{quote}崩漏（子宫异常出血）是妇科常见急症，不可掉以轻心。中医治疗崩漏遵循明代方约之提出的塞流、澄源、复旧三法：（1）塞流——急则治标，止血为先（如用三七、地榆炭）；（2）澄源——审因论治，清除病因（如血热者清热凉血，气虚者补气摄血）；（3）复旧——善后调理，重建正常月经周期（补肾调冲任）。西医根据病因分类（排卵障碍性、结构性等），治疗包括激素治疗、抗纤溶药物、宫腔镜手术等。异常出血需及时就医，排除器质性病变。\end{quote}
\noindent\textit{资料依据见本页脚注。}
\footnote{Mikes B A, Vadakekut E S, Sparzak P B. Abnormal uterine bleeding[M/OL]//StatPearls. Treasure Island: StatPearls Publishing, 2025[2026-06-18]. \href{https://www.ncbi.nlm.nih.gov/books/NBK532913/}{在线链接}.}
\footnote{National Institute for Health and Care Excellence. Heavy menstrual bleeding: assessment and management (NG88)[EB/OL]. 2018-03-14, last reviewed 2024-12-19[2026-06-18]. \href{https://www.nice.org.uk/guidance/ng88}{在线链接}.}
\footnote{世界中医药学会联合会, 中华中医药学会. 国际中医临床实践指南: 崩漏(2019-10-11)[J]. 世界中医药, 2021, 16(6): 870-877.}
\footnote{傅山. 傅青主女科: 血崩[M/OL]. 维基文库, [2026-06-18]. \href{https://zh.wikisource.org/zh-hans/傅青主女科}{在线链接}.}

\subsubsection{第六拍·濠州夺婚}\label{question:zhaomin:6}
\textbf{情节位置：}\hyperref[plot:zhaomin:6]{濠州夺婚完整剧情}。

\textbf{题干：}大喜大悲等剧烈情绪波动对女性健康影响深远。中医认为怒伤肝、喜伤心、思伤脾、悲伤肺、恐伤肾。以下关于女性情绪管理与健康的关系，哪一项不正确？

\textbf{选项与答案}
\begin{enumerate}[label=\Alph*.,leftmargin=*]
  \item 长期压力可能扰乱睡眠、月经节律，并加重经前期不适
  \item 中医逍遥散为疏肝解郁经典方，适用于肝郁血虚证
  \item 女性情绪波动只影响心理健康，与生殖系统无关\textbf{（正确答案）}
  \item 规律运动（如八段锦、瑜伽）和社交倾诉有助于调节情绪
\end{enumerate}
\textbf{判定依据、推导与科普边界}
\begin{quote}心理压力与身体健康并非两条互不相干的线。持续压力可通过下丘脑—垂体—肾上腺轴及下丘脑—垂体—卵巢轴影响睡眠、食欲、排卵和月经节律，也可能加重经前期障碍、围绝经期症状或焦虑抑郁。中医将部分表现辨为肝郁、心脾两虚等证候，逍遥散是疏肝养血的经典方之一，但必须由合格医师辨证，不能把所有情绪或月经问题都归结为肝郁并自行服药。规律运动、睡眠、社会支持与必要的心理/精神专科帮助，都是可靠的应对途径。\end{quote}
\noindent\textit{资料依据见本页脚注。}
\footnote{American College of Obstetricians and Gynecologists. Management of premenstrual disorders: ACOG Clinical Practice Guideline No. 7[J]. Obstetrics \& Gynecology, 2023, 142(6): 1516-1533. DOI：\nolinkurl{10.1097/AOG.0000000000005426}.}
\footnote{Reilly T J, Patel S, Unachukwu I C, et al. The prevalence of premenstrual dysphoric disorder: systematic review and meta-analysis[J]. Journal of Affective Disorders, 2024, 349: 534-540. DOI：\nolinkurl{10.1016/j.jad.2024.01.066}.}
\footnote{Chan K, Rubtsova A A, Clark C J. Exploring diagnosis and treatment of premenstrual dysphoric disorder in the U.S. healthcare system: a qualitative investigation[J]. BMC Women's Health, 2023, 23: 272. DOI：\nolinkurl{10.1186/s12905-023-02334-y}.}
\footnote{太医局. 太平惠民和剂局方: 逍遥散[M/OL]. 中医笈成, [2026-06-18]. \href{https://jicheng.tw/tcm/book/太平惠民和剂局方_1/index.html}{在线链接}.}

\subsubsection{第七拍·少室山下}\label{question:zhaomin:7}
\textbf{情节位置：}\hyperref[plot:zhaomin:7]{少室山下完整剧情}。

\textbf{题干：}张无忌提到脾胃为气血生化之源。中医认为女性补血先须健脾。以下关于女性营养性贫血的说法，哪一项是正确的？

\textbf{选项与答案}
\begin{enumerate}[label=\Alph*.,leftmargin=*]
  \item 缺铁性贫血只需多吃红枣、红糖即可
  \item 动物性铁（血红素铁）的吸收率高于植物性铁（非血红素铁）\textbf{（正确答案）}
  \item 贫血患者可自行服用铁剂，剂量越大越好
  \item 茶和咖啡有助于铁的吸收，贫血者可多饮
\end{enumerate}
\textbf{判定依据、推导与科普边界}
\begin{quote}动物性食物中的血红素铁通常比植物性非血红素铁更易吸收。红枣和红糖不能替代富含铁的均衡饮食或规范铁剂治疗；植物性铁可与富含维生素C的食物同食以帮助吸收。浓茶、咖啡中的多酚会影响非血红素铁吸收，可与正餐和铁剂错开。铁剂应在明确缺铁及病因后遵医嘱服用，过量可中毒。中医所称血虚不等同于实验室诊断的贫血，四物汤、归脾汤等方剂也须辨证处方，不能代替血常规、铁蛋白检查和病因治疗。\end{quote}
\noindent\textit{资料依据见本页脚注。}
\footnote{World Health Organization. WHO global anaemia estimates: key findings, 2025[R/OL]. Geneva: World Health Organization, 2025[2026-06-18]. \href{https://www.who.int/publications/i/item/9789240113930}{在线链接}.}
\footnote{中国营养学会. 中国居民膳食营养素参考摄入量(2023版)[M]. 北京: 人民卫生出版社, 2023.}
\footnote{National Institutes of Health Office of Dietary Supplements. Iron: fact sheet for health professionals[EB/OL]. 2025-09-18[2026-06-18]. \href{https://ods.od.nih.gov/factsheets/Iron-HealthProfessional/}{在线链接}.}
\footnote{DeLoughery T G, Jackson C S, Ko C W, et al. AGA clinical practice update on management of iron deficiency anemia: expert review[J]. Clinical Gastroenterology and Hepatology, 2024, 22(8): 1575-1583. DOI：\nolinkurl{10.1016/j.cgh.2024.03.046}.}

\subsubsection{第八拍·朔漠风霜}\label{question:zhaomin:8}
\textbf{情节位置：}\hyperref[plot:zhaomin:8]{朔漠风霜完整剧情}。

\textbf{题干：}赵敏提到古代女子以铅粉敷脸。关于女性皮肤健康与防晒，以下哪一项是正确的？

\textbf{选项与答案}
\begin{enumerate}[label=\Alph*.,leftmargin=*]
  \item 古代铅粉美白安全有效，现代也可使用
  \item 紫外线只导致皮肤晒黑，不会引起皮肤癌
  \item 涂抹防晒霜（SPF≥30）是预防皮肤光老化和皮肤癌的重要措施\textbf{（正确答案）}
  \item 阴天和冬天不需要防晒
\end{enumerate}
\textbf{判定依据、推导与科普边界}
\begin{quote}古代铅粉可能造成重金属暴露；现代化妆品法规禁止或严格限制相关有害物质及杂质限量，来源不明的速效美白产品尤其需要警惕。UVA与UVB会造成晒伤、光老化并增加皮肤癌风险，阴天和冬季也可能存在有效紫外线暴露。防护应结合遮阳、衣帽、避免正午暴晒与足量广谱防晒霜（通常建议SPF30及以上），并按出汗、游泳和产品说明补涂。传统草本外用方也可能刺激或致敏，内服中药更不应作为常规美容保健品自行使用。\end{quote}
\noindent\textit{资料依据见本页脚注。}
\footnote{World Health Organization. Radiation: protecting against skin cancer[EB/OL]. 2024-06-25[2026-06-18]. \href{https://www.who.int/news-room/questions-and-answers/item/radiation-protecting-against-skin-cancer}{在线链接}.}
\footnote{Centers for Disease Control and Prevention. Sun safety facts[EB/OL]. 2026-02-04[2026-06-18]. \href{https://www.cdc.gov/skin-cancer/sun-safety/index.html}{在线链接}.}
\footnote{吴谦, 等. 医宗金鉴·外科心法要诀: 黧黑皯〔黑曾〕[M/OL]. 维基文库, [2026-06-18]. \href{https://zh.wikisource.org/zh-hans/醫宗金鑒/外科卷上}{在线链接}.}
\footnote{Morriss S, Scardamaglia L. Sun protection: a practical guide for health professionals[J]. Australian Prescriber, 2025, 48(5): 173-178. DOI：\nolinkurl{10.18773/austprescr.2025.046}.}

\subsubsection{第十拍·父女恩仇}\label{question:zhaomin:10}
\textbf{情节位置：}\hyperref[plot:zhaomin:10]{父女恩仇完整剧情}。

\textbf{题干：}中年以后，无论男女都面临激素水平变化带来的身心挑战。以下关于围绝经期（更年期）女性健康的说法，哪一项是正确的？

\textbf{选项与答案}
\begin{enumerate}[label=\Alph*.,leftmargin=*]
  \item 更年期症状无法缓解，只能忍受
  \item 激素替代疗法（HRT）对所有更年期女性都适用，且无风险
  \item 围绝经期症状可以干预；绝经激素治疗对合适人群有效，但须由医师个体化评估获益、风险与禁忌证\textbf{（正确答案）}
  \item 更年期女性不需要进行妇科检查
\end{enumerate}
\textbf{判定依据、推导与科普边界}
\begin{quote}围绝经期是绝经前后的一段生理过渡，开始年龄和持续时间存在个体差异。潮热、盗汗、睡眠或泌尿生殖道症状均可干预。绝经激素治疗是缓解血管舒缩症状和泌尿生殖综合征的有效方法之一，但并非人人适用，需结合年龄、绝经年限、症状、血栓与肿瘤等风险个体化决策。中医可将相关表现辨为肾虚、肝郁等不同证候并辨证施治，不应自行套用成药。宫颈癌、乳腺癌、骨质疏松和心血管风险管理应按年龄、既往结果与个人风险执行，而非一律增加检查。\end{quote}
\noindent\textit{资料依据见本页脚注。}
\footnote{中国中西医结合学会妇产科专业委员会. 更年期综合征中西医结合诊治指南(2023年版)[J]. 中国实用妇科与产科杂志, 2023, 39(8): 799-808. DOI：\nolinkurl{10.19538/j.fk2023080109}.}
\footnote{中华医学会妇产科学分会绝经学组. 中国绝经管理与绝经激素治疗指南2023版[J]. 中华妇产科杂志, 2023, 58(1): 4-21. DOI：\nolinkurl{10.3760/cma.j.cn112141-20221118-00706}.}
\footnote{The North American Menopause Society. The 2022 hormone therapy position statement of The North American Menopause Society[J]. Menopause, 2022, 29(7): 767-794. DOI：\nolinkurl{10.1097/GME.0000000000002028}.}
\footnote{World Health Organization. Menopause[EB/OL]. 2024-10-16[2026-06-18]. \href{https://www.who.int/news-room/fact-sheets/detail/menopause}{在线链接}.}

\subsubsection{第十一拍·玄冥寒掌}\label{question:zhaomin:11}
\textbf{情节位置：}\hyperref[plot:zhaomin:11]{玄冥寒掌完整剧情}。

\textbf{题干：}中医八纲辨证（阴阳、表里、寒热、虚实）是诊断的基础。关于寒证与热证在女性健康中的表现，以下哪一项是正确的？

\textbf{选项与答案}
\begin{enumerate}[label=\Alph*.,leftmargin=*]
  \item 痛经只有寒证一种类型，只需多喝热水
  \item 寒证痛经（小腹冷痛、喜温喜按）与热证痛经（灼热痛、经色深红）治法不同，需辨证施治\textbf{（正确答案）}
  \item 寒热辨证对女性疾病没有参考价值
  \item 所有妇科疾病都用热性药治疗
\end{enumerate}
\textbf{判定依据、推导与科普边界}
\begin{quote}妇科疾病中寒热辨证至关重要，治法截然相反。寒证痛经（宫寒型）：小腹冷痛、得热则减、经色暗有块、畏寒肢冷，治以温经散寒；热证痛经（湿热/郁热型）：小腹灼痛、经色深红质稠、口干舌红，治以清热凉血。若寒热不分，如宫寒者误用寒凉药则雪上加霜。现代医学认为原发性痛经常与前列腺素升高、子宫收缩和缺血相关；继发性痛经需警惕子宫内膜异位症、盆腔炎等。热疗对原发性痛经有一定缓痛证据；艾灸、针刺等中医外治应辨证使用，并由专业人员指导。\end{quote}
\noindent\textit{资料依据见本页脚注。}
\footnote{佚名. 黄帝内经·素问·至真要大论[M/OL]. 中国哲学书电子化计划, [2026-06-18]. \href{https://ctext.org/huangdi-neijing/zhi-zhen-yao-da-lun/zh}{在线链接}.}
\footnote{Kirsch E, Rahman S, Kerolus K, et al. Dysmenorrhea, a narrative review of therapeutic options[J]. Journal of Pain Research, 2024, 17: 2657-2666. DOI：\nolinkurl{10.2147/JPR.S459584}.}
\footnote{Yuan D, Liu Y, Chen Z, et al. Heat therapy for primary dysmenorrhea: a systematic review and meta-analysis[J]. Frontiers in Medicine, 2026, 12: 1730505. DOI：\nolinkurl{10.3389/fmed.2025.1730505}.}
\footnote{傅山. 傅青主女科: 经水将来脐下先疼痛[M/OL]. 维基文库, [2026-06-18]. \href{https://zh.wikisource.org/zh-hans/傅青主女科}{在线链接}.}

\subsubsection{第十二拍·大都宫闱}\label{question:zhaomin:12}
\textbf{情节位置：}\hyperref[plot:zhaomin:12]{大都宫闱完整剧情}。

\textbf{题干：}御医所言乌头有剧毒，炮制得法可为良药，体现了中医以毒攻毒的思想。以下关于中药毒性认识的说法，哪一项是不正确的？

\textbf{选项与答案}
\begin{enumerate}[label=\Alph*.,leftmargin=*]
  \item 附子（乌头子根）经炮制后可大幅降低乌头碱含量
  \item 朱砂（硫化汞）在中医中曾用于安神，但因含汞现已严格控制使用
  \item 所有中草药都无毒副作用，可长期随意服用\textbf{（正确答案）}
  \item 马钱子含士的宁（番木鳖碱），过量可致惊厥死亡
\end{enumerate}
\textbf{判定依据、推导与科普边界}
\begin{quote}天然不等于安全。乌头类可诱发恶性心律失常，马钱子中的士的宁可致惊厥，朱砂和雄黄涉及汞、砷暴露风险；这些药材均有严格的炮制、剂量、疗程和适应证要求。安全用药要看具体品种、来源、炮制、配伍、剂量、基础疾病和合并用药，不能笼统宣称某一中成药常规剂量下都安全。即使是常见中成药，也可能有禁忌、不良反应或药物相互作用，应按说明书并在医师或药师指导下使用。\end{quote}
\noindent\textit{资料依据见本页脚注。}
\footnote{佚名. 神农本草经[M/OL]. 中医笈成, [2026-06-18]. \href{https://jicheng.tw/tcm/book/神農本草經_2/index.html}{在线链接}.}
\footnote{国家药品监督管理局, 国家卫生健康委员会. 关于颁布2025年版《中华人民共和国药典》的公告(2025年第29号)[EB/OL]. 2025-03-25[2026-06-18]. \href{https://www.nmpa.gov.cn/xxgk/fgwj/gzwj/gzwjyp/20250325183810122.html}{在线链接}.}
\footnote{Zhou C, Luo S, Tang J, et al. Poisoning associated with consumption of a homemade medicinal liquor - Chongqing, China, 2018[J]. MMWR. Morbidity and Mortality Weekly Report, 2022, 71(16): 569-573. DOI：\nolinkurl{10.15585/mmwr.mm7116a2}.}
\footnote{Lawson C, McCabe D J, Feldman R. A narrative review of aconite poisoning and management[J]. Journal of Intensive Care Medicine, 2025, 40(8): 811-817. DOI：\nolinkurl{10.1177/08850666241245703}.}

\subsubsection{第十四拍·光明顶上}\label{question:zhaomin:14}
\textbf{情节位置：}\hyperref[plot:zhaomin:14]{光明顶上完整剧情}。

\textbf{题干：}女性运动不当可导致运动损伤。以下关于女性运动与健康的说法，哪一项是正确的？

\textbf{选项与答案}
\begin{enumerate}[label=\Alph*.,leftmargin=*]
  \item 女性经期绝对不能运动
  \item 凯格尔运动（盆底肌训练）对女性盆底健康尤其重要，可辅助防治尿失禁等盆底功能障碍\textbf{（正确答案）}
  \item 运动量越大对健康越有利
  \item 女性练力量运动会变得太壮，应避免
\end{enumerate}
\textbf{判定依据、推导与科普边界}
\begin{quote}盆底肌训练是女性尿失禁等盆底功能障碍的一线保守治疗之一，但关键在于正确识别、收缩并充分放松盆底肌；有疼痛、排尿困难或盆底过度紧张者宜先接受专业评估。月经本身不是停止运动的理由，可根据疼痛、出血量和个人感受调整强度。运动应循序渐进；长期能量摄入不足与过量训练可造成相对能量缺乏和功能性下丘脑性闭经。适当抗阻训练有助于肌力与骨骼健康，八段锦、太极拳等也可作为低至中等强度活动选择。\end{quote}
\noindent\textit{资料依据见本页脚注。}
\footnote{Woodley S J, Boyle R, Cody J D, et al. Pelvic floor muscle training for urinary incontinence in women[J]. Cochrane Database of Systematic Reviews, 2024, 12:CD009508. PMID:39704322.}
\footnote{American College of Obstetricians and Gynecologists. Urinary incontinence: Kegel exercises[EB/OL]. \href{https://www.acog.org/womens-health/faqs/urinary-incontinence}{在线链接}.}
\footnote{National Institute for Health and Care Excellence. Urinary incontinence and pelvic organ prolapse in women: management (NG123)[S]. 2019 (updated). \href{https://www.nice.org.uk/guidance/ng123}{在线链接}.}
\footnote{Mountjoy M, Sundgot-Borgen J, Burke L, et al. 2023 international Olympic Committee consensus statement on relative energy deficiency in sport (REDs)[J]. British Journal of Sports Medicine, 2023, 57(17):1073–1098. DOI：\nolinkurl{10.1136/bjsports-2023-106994}.}

\subsubsection{第十五拍·波斯异术}\label{question:zhaomin:15}
\textbf{情节位置：}\hyperref[plot:zhaomin:15]{波斯异术完整剧情}。

\textbf{题干：}波斯医学中常用的藏红花（番红花）在中医和现代医学中均有重要地位。以下关于藏红花的说法，哪一项是正确的？

\textbf{选项与答案}
\begin{enumerate}[label=\Alph*.,leftmargin=*]
  \item 藏红花与普通红花（草红花）完全相同
  \item 藏红花可用于治疗经闭、痛经、产后瘀阻，但孕妇禁用\textbf{（正确答案）}
  \item 藏红花价格便宜，可大量服用
  \item 藏红花只在中国有产
\end{enumerate}
\textbf{判定依据、推导与科普边界}
\begin{quote}藏红花（Crocus sativus，鸢尾科）与草红花（Carthamus tinctorius，菊科）是两种不同植物。藏红花取雌蕊柱头入药，产量极低（约15万朵花产1kg干柱头），故价格昂贵（红色黄金）。中医认为藏红花性甘微寒，归心肝经，功效活血化瘀、凉血解毒、解郁安神。现代研究证实其活性成分藏红花素（crocin）和藏红花醛（safranal）具有：（1）抗抑郁作用（多项RCT证实对轻中度抑郁有效）；（2）缓解经前期综合征（PMS）；（3）改善原发性痛经。因有活血化瘀和兴奋子宫作用，孕妇禁用。主要产自伊朗（>90\%）、西班牙、印度克什米尔，中国西藏有少量种植。\end{quote}
\noindent\textit{资料依据见本页脚注。}
\footnote{国家药品监督管理局, 国家卫生健康委员会. 关于颁布2025年版《中华人民共和国药典》的公告（2025年第29号）[S]. 2025-03-20. \href{https://www.nmpa.gov.cn/xxgk/fgwj/gzwj/gzwjyp/20250325183810122.html}{在线链接}.}
\footnote{Mohammadi M M, Karimi Z. Effect of saffron on premenstrual syndrome and dysmenorrhea: a systematic review and meta-analysis[J]. Korean Journal of Family Medicine, 2026, 47(1):69–80. DOI：\nolinkurl{10.4082/kjfm.24.0259}.}
\footnote{Hasheminasab F S, Azimi M, Raeiszadeh M, et al. Therapeutic effects of saffron (Crocus sativus L.) on female reproductive system disorders: a systematic review[J]. Phytotherapy Research, 2024. PMID:38558480.}
\footnote{卫生福利部国家中医药研究所. 番红花与红花[EB/OL]. \href{https://www.nricm.edu.tw/p/404-1000-5146.php?Lang=zh-tw}{在线链接}.}

\subsubsection{第十六拍·襄阳遗梦}\label{question:zhaomin:16}
\textbf{情节位置：}\hyperref[plot:zhaomin:16]{襄阳遗梦完整剧情}。

\textbf{题干：}赵敏提到战地急救的基本技能。日常生活中，女性也需掌握一些急救知识，尤其是面对妇女儿童常见意外的处理。以下关于家庭急救的说法，哪一项是错误的？

\textbf{选项与答案}
\begin{enumerate}[label=\Alph*.,leftmargin=*]
  \item 鼻出血时应仰头，并把纸团深塞入鼻腔\textbf{（正确答案）}
  \item 烫伤后应立即用流动清水冲洗至少20分钟
  \item 心肺复苏（CPR）按压深度为5-6厘米，频率100-120次/分
  \item 海姆立克急救法用于气道异物梗阻
\end{enumerate}
\textbf{判定依据、推导与科普边界}
\begin{quote}鼻出血时应坐直、身体略前倾，用拇指和食指持续捏紧鼻翼柔软处10-15分钟并用口呼吸；不要仰头，也不要反复松手查看或把纸团深塞入鼻腔。若持续压迫后仍大量出血、因外伤引起、伴头晕乏力，或正在使用抗凝药，应及时就医。成人心肺复苏胸外按压通常为每分钟100-120次、深度至少5厘米并避免超过6厘米；气道异物处理应按患者年龄、是否清醒及能否有效咳嗽选择拍背、腹部冲击或胸部冲击。\end{quote}
\noindent\textit{资料依据见本页脚注。}
\footnote{Panchal A R, Bartos J A, Cabanas J G, et al. 2024 American Heart Association and American Red Cross guidelines for first aid[J]. Circulation, 2024. DOI：\nolinkurl{10.1161/CIR.0000000000001281}.}
\footnote{American Heart Association. 2025 American Heart Association guidelines for CPR and emergency cardiovascular care[J]. Circulation, 2025, 152(suppl 2). DOI：\nolinkurl{10.1161/CIR.0000000000001372}.}
\footnote{Mayo Clinic Staff. Nosebleeds: first aid[EB/OL]. \href{https://www.mayoclinic.org/first-aid/first-aid-nosebleeds/basics/art-20056683}{在线链接}.}
\footnote{American Red Cross. Nosebleed first aid guidelines[EB/OL]. \href{https://www.redcross.org/take-a-class/resources/guidelines}{在线链接}.}

\subsubsection{第十八拍·天下归心}\label{question:zhaomin:18}
\textbf{情节位置：}\hyperref[plot:zhaomin:18]{天下归心完整剧情}。

\textbf{题干：}赵敏希望将十八关中所学医道整理成书，帮助天下女子。以下关于女性全生命周期健康管理的说法，哪一项最全面地概括了女性健康的核心要点？

\textbf{选项与答案}
\begin{enumerate}[label=\Alph*.,leftmargin=*]
  \item 女性只需要关注生育期的健康
  \item 女性全生命周期健康管理应从青春期开始，涵盖育龄期、围绝经期和老年期，包括生理、心理和社会适应三个维度\textbf{（正确答案）}
  \item 女性健康管理只需要做妇科检查
  \item 女性过了生育年龄就不需要特别关注健康了
\end{enumerate}
\textbf{判定依据、推导与科普边界}
\begin{quote}世界卫生组织（WHO）强调以全生命周期（life-course approach）理解健康：青春期、育龄期、围绝经期与老年期彼此相连，不应只在生病或备孕时才关注身体。各阶段重点包括：青春期的月经与营养教育、按国家免疫程序接种HPV疫苗；育龄期的避孕或孕前保健，以及按年龄与风险开展宫颈癌和乳腺癌筛查；围绝经期的症状管理、心血管和骨折风险评估；老年期的防跌倒、认知、盆底与慢病管理。具体筛查起始年龄和间隔须遵循所在地最新指南与个人风险，不能用一套固定检查覆盖所有女性。中医治未病思想可与现代预防医学对话，但预防仍须建立在疫苗、筛查、生活方式和规范诊疗之上。\end{quote}
\noindent\textit{资料依据见本页脚注。}
\footnote{World Health Organization. Global strategy for women’s, children’s and adolescents’ health (2016–2030)［EB/OL］. \href{https://platform.who.int/data/maternal-newborn-child-adolescent-ageing/global-strategy-data}{在线链接}.}
\footnote{World Health Organization. Women’s health［EB/OL］. \href{https://www.who.int/health-topics/women-s-health}{在线链接}.}
\footnote{Pan American Health Organization. Healthy life course［EB/OL］. \href{https://www.paho.org/en/topics/healthy-life-course}{在线链接}.}
\footnote{国务院. 中国妇女发展纲要（2021—2030年）［Z］. 国发〔2021〕16号［EB/OL］. 2021-09-08. \href{https://www.mee.gov.cn/zcwj/gwywj/202109/t20210927_954260.shtml}{在线链接}.}
\footnote{佚名. 黄帝内经·素问·四气调神大论［M］. 北京：中国哲学书电子化计划（CTEXT）数据库［EB/OL］. \href{https://ctext.org/huangdi-neijing/si-qi-diao-shen-da-lun/zh}{在线链接}.}

\subsection{周芷若线：汉水白莲}

\subsubsection{第一拍·汉水初遇}\label{question:zhouzhiruo:1}
\textbf{情节位置：}\hyperref[plot:zhouzhiruo:1]{汉水初遇完整剧情}。

\textbf{题干：}周芷若父亲死于长期营养不良。在战乱年代，女性因生理特点更易发生缺铁性贫血。以下关于女性贫血的说法，哪一项是正确的？

\textbf{选项与答案}
\begin{enumerate}[label=\Alph*.,leftmargin=*]
  \item 贫血就是血少，多喝水就可以了
  \item 月经过多、妊娠、素食饮食均可导致或加重缺铁性贫血\textbf{（正确答案）}
  \item 贫血只需吃素就能改善
  \item 女性面色苍白是气血充足的正常表现
\end{enumerate}
\textbf{判定依据、推导与科普边界}
\begin{quote}女性因月经失血、妊娠需求增加、产后失血等因素，是缺铁性贫血的高发人群。WHO 2025版全球贫血估计显示，2023年15-49岁女性贫血患病率约30.7\%。常见原因：（1）月经量过多（>80ml/周期）；（2）含铁食物摄入不足（素食者尤为常见，因植物铁吸收率低）；（3）妊娠期血容量增加、胎儿需求；（4）慢性失血（如消化道出血）。中医属血虚范畴，常见证型有：心血虚（心悸失眠，归脾汤）、肝血虚（头晕目眩，四物汤）、气血两虚（面色萎黄，八珍汤）。诊断需查血常规+铁蛋白，治疗应针对病因+补充铁剂（遵医嘱）。\end{quote}
\noindent\textit{资料依据见本页脚注。}
\footnote{World Health Organization. WHO global anaemia estimates: key findings, 2025[R]. Geneva: World Health Organization, 2025. Available: \href{https://www.who.int/publications/i/item/9789240113930}{在线链接}.}
\footnote{Mayo Clinic. Iron deficiency anemia: Symptoms and causes[EB/OL]. Rochester: Mayo Foundation for Medical Education and Research, 2024[2026-06-18]. \href{https://www.mayoclinic.org/diseases-conditions/iron-deficiency-anemia/symptoms-causes/syc-20355034}{在线链接}}
\footnote{Pavord S, Daru J, Prasannan N, et al. UK guidelines on the management of iron deficiency in pregnancy[J]. British Journal of Haematology, 2020, 188(6): 819-830. DOI：\nolinkurl{10.1111/bjh.16221}.}
\footnote{中国营养学会. 中国居民膳食营养素参考摄入量（2023版）[M]. 北京: 人民卫生出版社, 2023. ISBN: 978-7-117-35080-8.}

\subsubsection{第二拍·峨眉入门}\label{question:zhouzhiruo:2}
\textbf{情节位置：}\hyperref[plot:zhouzhiruo:2]{峨眉入门完整剧情}。

\textbf{题干：}灭绝师太所言的天癸，在中医理论中指促进生殖功能成熟的精微物质。关于女性青春期发育和月经初潮，以下哪一项不正确？

\textbf{选项与答案}
\begin{enumerate}[label=\Alph*.,leftmargin=*]
  \item 月经初潮是女性青春期的标志，通常发生在10-16岁
  \item 青春期女性应完全避免运动，以免损伤天癸\textbf{（正确答案）}
  \item 经期适度运动（如瑜伽、散步）有助于缓解经期不适
  \item 青春期营养不足可导致初潮延迟或月经稀发
\end{enumerate}
\textbf{判定依据、推导与科普边界}
\begin{quote}适度的运动不会损伤天癸，反而有助于青春期发育。正确的做法是：（1）经期避免剧烈运动和重体力劳动（中医忌过劳），但适度活动如散步、瑜伽、八段锦有助于气血运行，缓解痛经；（2）青春期需要充足营养——蛋白质、铁、钙、锌等，节食或营养不良可致初潮延迟（<16岁称原发性闭经）或继发性闭经；（3）初潮后前1-2年月经可能不规律，属正常现象（无排卵性月经），但若持续不规律超过3年需就医；（4）青春期女孩需进行月经卫生教育、HPV疫苗接种（9-14岁最佳）。\end{quote}
\noindent\textit{资料依据见本页脚注。}
\footnote{佚名. 黄帝内经·素问·上古天真论［M］. 北京：中国哲学书电子化计划（CTEXT）数据库，\href{https://ctext.org/huangdi-neijing/shang-gu-tian-zhen-lun/zh}{在线链接}.}
\footnote{American College of Obstetricians and Gynecologists. Menstruation in girls and adolescents: using the menstrual cycle as a vital sign［J］. Obstetrics \& Gynecology, 2015, 126(6): e143–e146. DOI：\nolinkurl{10.1097/AOG.0000000000001215}.}
\footnote{World Health Organization. Human papillomavirus vaccines: WHO position paper, December 2022［J］. Weekly Epidemiological Record, 2022, 97(50): 645–672. \href{https://www.who.int/publications/i/item/who-wer9750-645-672}{在线链接}.}
\footnote{World Health Organization. Adolescent health［EB/OL］. \href{https://www.who.int/health-topics/adolescent-health/}{在线链接} (accessed 2026-06-18).}

\subsubsection{第三拍·光明一剑}\label{question:zhouzhiruo:3}
\textbf{情节位置：}\hyperref[plot:zhouzhiruo:3]{光明一剑完整剧情}。

\textbf{题干：}周芷若在极端压力下刺伤张无忌——这一心理创伤影响了她此后的人生轨迹。关于创伤后应激与女性心理健康，以下哪一项是正确的？

\textbf{选项与答案}
\begin{enumerate}[label=\Alph*.,leftmargin=*]
  \item 心理创伤只是暂时的，时间可以治愈一切，不需要任何干预
  \item 女性因生理和社会因素，创伤后应激障碍（PTSD）发病率约为男性的2倍\textbf{（正确答案）}
  \item 中医认为情绪与身体健康无关
  \item 创伤后应激反应越强烈说明心理素质越差
\end{enumerate}
\textbf{判定依据、推导与科普边界}
\begin{quote}流行病学研究通常发现女性发生PTSD的风险高于男性，约可达到两倍，但具体比例随人群、创伤类型和研究方法而变；差异与创伤暴露类型、社会环境和生物学因素共同相关，不能归咎于心理素质差。PTSD核心症状包括闯入性再体验、回避、认知和情绪的负性改变以及警觉性增高。循证治疗以创伤聚焦心理治疗（如创伤聚焦CBT、EMDR等）为重要选择，部分患者需要规范药物治疗。中医可从惊悸、郁证等角度辨证辅助，但不能用固定方剂替代创伤评估与专业治疗；如有自伤念头或无法保证安全，应立即求助。\end{quote}
\noindent\textit{资料依据见本页脚注。}
\footnote{American Psychiatric Association. Diagnostic and statistical manual of mental disorders: DSM-5-TR[M]. Washington, DC: American Psychiatric Publishing, 2022. DOI：\nolinkurl{10.1176/appi.books.9780890425787}.}
\footnote{U.S. Department of Veterans Affairs. PTSD: DSM-5 criteria[EB/OL]. \href{https://www.ptsd.va.gov/professional/treat/essentials/dsm5_ptsd.asp}{在线链接}.}
\footnote{Hiscox LV, Sharp TH, Olff M, Seedat S, Halligan SL. Sex-based contributors to and consequences of post-traumatic stress disorder[J]. Current Psychiatry Reports, 2023, 25:233–245. DOI：\nolinkurl{10.1007/s11920-023-01421-z}.}
\footnote{Merck Manual Professional Edition. Posttraumatic stress disorder (PTSD)[EB/OL]. \href{https://www.merckmanuals.com/professional/psychiatric-disorders/anxiety-and-trauma-and-stressor-related-disorders/posttraumatic-stress-disorder-ptsd}{在线链接}. (accessed 2026-06-18)}

\subsubsection{第五拍·灵蛇毒计}\label{question:zhouzhiruo:5}
\textbf{情节位置：}\hyperref[plot:zhouzhiruo:5]{灵蛇毒计完整剧情}。

\textbf{题干：}周芷若使用的十香软筋散虽为虚构，但药物安全是真实世界的重要议题。以下关于药物安全与毒理学的说法，哪一项是正确的？

\textbf{选项与答案}
\begin{enumerate}[label=\Alph*.,leftmargin=*]
  \item 中药都是天然的，不会有任何毒副作用
  \item 附子、乌头等中药含乌头碱，过量可导致心律失常甚至死亡\textbf{（正确答案）}
  \item 任何药物都可以随意加大剂量以增强疗效
  \item 药物毒性与剂量无关
\end{enumerate}
\textbf{判定依据、推导与科普边界}
\begin{quote}天然不等于安全。乌头类所含双酯型二萜生物碱可持续激活电压门控钠通道，引起感觉异常、低血压和恶性心律失常，严重时致死。规范炮制与煎煮可降低毒性，但不能把有毒药材变成可随意加量的保健品；具体降毒幅度受品种、工艺和检测方法影响，不宜用单一百分比概括。药物风险与剂量、炮制、配伍、基础疾病及合并用药有关。对乙酰氨基酚等西药超量同样可造成严重中毒，所有药物都应按说明或医嘱使用。\end{quote}
\noindent\textit{资料依据见本页脚注。}
\footnote{中国哲学书电子化计划. 神农本草经[EB/OL]. \href{https://jicheng.tw/tcm/book/神農本草經_2/index.html}{在线链接}.}
\footnote{国家药品监督管理局, 国家卫生健康委员会. 关于颁布2025年版《中华人民共和国药典》的公告（2025年第29号）[S]. 2025-03-20. \href{https://www.nmpa.gov.cn/xxgk/fgwj/gzwj/gzwjyp/20250325183810122.html}{在线链接}.}
\footnote{Centers for Disease Control and Prevention. Poisoning associated with consumption of a homemade medicinal drink — California, 2022[J/OL]. Morbidity and Mortality Weekly Report, 2022, 71(16). \href{https://www.cdc.gov/mmwr/volumes/71/wr/mm7116a2.htm}{在线链接}.}
\footnote{Lawson C, McCabe DJ, Feldman R. A narrative review of aconite poisoning and management[J]. Journal of Intensive Care Medicine, 2025, 40(8):811–817. PMID:38613376.}

\subsubsection{第六拍·裂裳之痛}\label{question:zhouzhiruo:6}
\textbf{情节位置：}\hyperref[plot:zhouzhiruo:6]{裂裳之痛完整剧情}。

\textbf{题干：}失恋、离异等重大情感创伤对女性身心健康影响深远。以下关于情感与女性健康的说法，哪一项是正确的？

\textbf{选项与答案}
\begin{enumerate}[label=\Alph*.,leftmargin=*]
  \item 被背叛的一方不值得同情，是她们自己不够好
  \item 中医认为大喜大悲伤心伤肝，可致气机逆乱、月经失调\textbf{（正确答案）}
  \item 经历情感创伤后，应立刻投入新恋情来治愈
  \item 情感创伤只影响心理，对生理没有实质影响
\end{enumerate}
\textbf{判定依据、推导与科普边界}
\begin{quote}中医经典《黄帝内经》明确指出：怒伤肝、喜伤心、思伤脾、悲伤肺、恐伤肾。大喜（如婚礼）大悲（如被弃）的剧烈转换，最能扰乱气机——惊则气乱悲则气消。心主血脉，心伤则血行失常；肝主疏泄，肝郁则月经不调。现代医学证实：重大情感创伤可激活HPA轴（下丘脑-垂体-肾上腺轴），导致皮质醇持续升高，进而抑制HPG轴（下丘脑-垂体-性腺轴），引起排卵障碍、月经紊乱。心碎综合征（Takotsubo心肌病）更是情感打击导致心肌功能障碍的极端例子。正确的疗愈方式：给时间、寻求社交支持、专业心理咨询，而非仓促进入新关系。\end{quote}
\noindent\textit{资料依据见本页脚注。}
\footnote{中国哲学书电子化计划. 黄帝内经·素问·阴阳应象大论[EB/OL]. \href{https://ctext.org/huangdi-neijing/yin-yang-ying-xiang-da-lun/zhs}{在线链接}.}
\footnote{Wittstein I S, Thiemann D R, Lima J A C, et al. Neurohumoral features of myocardial stunning due to sudden emotional stress[J]. New England Journal of Medicine, 2005, 352(6):539–548. DOI：\nolinkurl{10.1056/NEJMoa043046}.}
\footnote{Sbarra D A, Whisman M A. Divorce, health, and socioeconomic status: An agenda for psychological science[J]. Current Opinion in Psychology, 2022, 43:75–78. DOI：\nolinkurl{10.1016/j.copsyc.2021.06.007}.}
\footnote{StatPearls Publishing. Takotsubo cardiomyopathy[EB/OL]. National Library of Medicine, 2026. \href{https://www.ncbi.nlm.nih.gov/books/NBK430798/}{在线链接}.}

\subsubsection{第七拍·九阴白骨}\label{question:zhouzhiruo:7}
\textbf{情节位置：}\hyperref[plot:zhouzhiruo:7]{九阴白骨完整剧情}。

\textbf{题干：}周芷若的指关节疼痛可能与过度使用导致的腱鞘炎、关节炎有关。以下关于女性关节健康的说法，哪一项是正确的？

\textbf{选项与答案}
\begin{enumerate}[label=\Alph*.,leftmargin=*]
  \item 关节疼痛是小问题，忍一忍就过去了
  \item 女性患类风湿关节炎和骨关节炎的风险高于男性，需重视早期症状\textbf{（正确答案）}
  \item 运动对关节只有伤害没有益处
  \item 关节疼痛只需热敷即可，无需就医
\end{enumerate}
\textbf{判定依据、推导与科普边界}
\begin{quote}女性关节疾病负担显著高于男性：（1）类风湿关节炎（RA）——女性患病率为男性的2-3倍，与雌激素对免疫系统的调节有关，好发于30-50岁；（2）骨关节炎（OA）——尤其在绝经后，女性手部、膝部OA发病率显著上升（雌激素对软骨有保护作用）；（3）狭窄性腱鞘炎（妈妈手、De Quervain病）在女性中也更常见。早期症状：晨僵（>30分钟需警惕RA）、关节肿痛、活动受限。防治：（1）适度运动（游泳、太极拳）对关节有保护作用，而非伤害；（2）保持健康体重（减重1kg可减少膝OA风险4\%）；（3）中医以通络止痛为法（独活寄生汤、针灸温针）；（4）及时就医，RA早期干预可改变疾病进程（Window of Opportunity概念）。\end{quote}
\noindent\textit{资料依据见本页脚注。}
\footnote{Di Matteo A, Bathon J M, Emery P. Rheumatoid arthritis［J］. The Lancet, 2023, 402(10416): 2019-2033. DOI：\nolinkurl{10.1016/S0140-6736(23)01525-8}.}
\footnote{Hunter D J, Bierma-Zeinstra S. Osteoarthritis［J］. The Lancet, 2019, 393(10182): 1745-1759. DOI：\nolinkurl{10.1016/S0140-6736(19)30417-9}.}
\footnote{Smolen J S, Aletaha D, McInnes I B. Rheumatoid arthritis［J］. Nature Reviews Disease Primers, 2018, 4: 18001. DOI：\nolinkurl{10.1038/nrdp.2018.1}.}
\footnote{中华中医药学会骨伤科分会. 膝骨关节炎中医临床诊疗指南［J］. 康复学报, 2019, 29(3): 1-7.}

\subsubsection{第八拍·峨眉掌门}\label{question:zhouzhiruo:8}
\textbf{情节位置：}\hyperref[plot:zhouzhiruo:8]{峨眉掌门完整剧情}。

\textbf{题干：}中年职业女性——如同周芷若执掌峨眉——常面临职场压力、家庭责任与自我健康管理的三重挑战。以下关于职业女性健康管理的说法，哪一项是合理的？

\textbf{选项与答案}
\begin{enumerate}[label=\Alph*.,leftmargin=*]
  \item 为了事业成功牺牲健康是值得的
  \item 合理安排作息、定期体检、保持社交支持是职业女性健康管理的关键\textbf{（正确答案）}
  \item 女性领导应该像男性一样强硬，不能示弱
  \item 熬夜和快餐是高效工作的合理代价
\end{enumerate}
\textbf{判定依据、推导与科普边界}
\begin{quote}职业女性健康管理应避免一张体检套餐包打天下：宫颈癌、乳腺癌、血压、血脂、血糖等筛查应按年龄、既往结果和个人风险安排，甲状腺或盆腔超声并非所有无症状女性每年必做。多数成人需要规律而充足的睡眠、每周至少150分钟中等强度活动并配合肌力训练；轮班工作造成的昼夜节律紊乱被IARC列为很可能对人类致癌，更应重视排班、睡眠与职业健康保护。稳定的社会支持、合理边界和及时求助同样重要；中药调理需辨证，不能替代循证筛查和心理支持。\end{quote}
\noindent\textit{资料依据见本页脚注。}
\footnote{International Agency for Research on Cancer. IARC monographs on the identification of carcinogenic hazards to humans. Volume 124: Night shift work［M］. Lyon: IARC, 2020. Available at: \href{https://publications.iarc.who.int/Book-And-Report-Series/Iarc-Monographs-On-The-Identification-Of-Carcinogenic-Hazards-To-Humans/Night-Shift-Work-2020}{在线链接}.}
\footnote{Shanafelt T D, West C P, Dyrbye L N, et al. Changes in burnout and satisfaction with work-life integration in physicians during the first 2 years of the COVID-19 pandemic［J］. Mayo Clinic Proceedings, 2022, 97(12): 2248-2258. DOI：\nolinkurl{10.1016/j.mayocp.2022.09.002}.}
\footnote{World Health Organization. Burn-out an occupational phenomenon: ICD-11［EB/OL］. \href{https://www.who.int/standards/classifications/frequently-asked-questions/burn-out-an-occupational-phenomenon}{在线链接}.}

\subsubsection{第十拍·黄衫女子}\label{question:zhouzhiruo:10}
\textbf{情节位置：}\hyperref[plot:zhouzhiruo:10]{黄衫女子完整剧情}。

\textbf{题干：}黄衫女子的话触及了一个深刻的问题——女性的自我价值与独立人格。研究表明，积极的女性榜样（role model）对女性心理健康有重要影响。以下关于女性自我价值与心理健康的说法，哪一项是正确的？

\textbf{选项与答案}
\begin{enumerate}[label=\Alph*.,leftmargin=*]
  \item 女性的价值取决于婚姻和家庭
  \item 建立独立的自我认同、发展多元社会角色（而非仅局限于家庭）有利于女性心理健康\textbf{（正确答案）}
  \item 独立女性不需要社交关系
  \item 女性应该完全效仿男性才能成功
\end{enumerate}
\textbf{判定依据、推导与科普边界}
\begin{quote}心理学家埃里克森（Erikson）提出自我同一性概念——健康的个体需要建立独立的、整合的自我认同。研究证实：女性拥有多元社会角色（职业、兴趣、社交等）与更好的心理健康相关。关键要素：（1）经济独立——是女性自主权的基础；（2）社会支持网络——不是不需要关系，而是拥有健康互惠的关系（非依附性关系）；（3）终身学习与成长——不断拓展自我边界。中医角度：独立不等于孤独——独阴不生，独阳不长，健康人格需要阴阳调和——既有独立的自我，也有连接他人的能力。古代女科医书将女性病证独立讨论，提醒医者认真看见女性身心处境；黄衫女子独立于终南山、小龙女独立于古墓，皆是金庸笔下独立的女性力量的象征。\end{quote}
\noindent\textit{资料依据见本页脚注。}
\footnote{Erikson E H. Identity and the life cycle［M］. New York: International Universities Press, 1959. Available at: \href{https://archive.org/details/psychologicaliss0001erik}{在线链接}.}
\footnote{World Health Organization. Women’s health［EB/OL］. \href{https://www.who.int/health-topics/women-s-health}{在线链接}.}
\footnote{National Academies of Sciences, Engineering, and Medicine. Advancing research on chronic conditions in women［M］. Washington, DC: The National Academies Press, 2024. DOI：\nolinkurl{10.17226/27757}.}
\footnote{萧壎. 女科经纶［M］. 清代刻本. 电子整理本: 中国哲学书电子化计划（CTEXT）［EB/OL］. \href{https://zh.wikisource.org/zh-hans/女科經綸/序}{在线链接}.}

\subsubsection{第十一拍·问心有愧}\label{question:zhouzhiruo:11}
\textbf{情节位置：}\hyperref[plot:zhouzhiruo:11]{问心有愧完整剧情}。

\textbf{题干：}脏躁是中医特有的概念，与现代医学中的抑郁症、焦虑症有相似之处。张仲景《金匮要略》以甘麦大枣汤治妇人脏躁。以下关于脏躁与女性情志疾病的说法，哪一项是正确的？

\textbf{选项与答案}
\begin{enumerate}[label=\Alph*.,leftmargin=*]
  \item 脏躁只是女性情绪化，不需要认真对待
  \item 围产期抑郁（产前/产后抑郁）是孕产期女性需重点筛查和支持的心理健康问题，发病率约10-15\%\textbf{（正确答案）}
  \item 中医的情志疾病与现代精神医学没有可对话之处
  \item 治疗脏躁只需吃药，心理疏导没有用
\end{enumerate}
\textbf{判定依据、推导与科普边界}
\begin{quote}脏躁首见于张仲景《金匮要略·妇人杂病脉证并治》：妇人脏躁，悲伤欲哭，象如神灵所作，数欠伸，甘麦大枣汤主之。此方仅甘草、小麦、大枣三味，却成千古名方——体现了中医简、便、验、廉的特点。围产期抑郁可发生在妊娠期至产后一年内，症状包括持续情绪低落、兴趣丧失、自责、焦虑，严重时可有伤害自己或婴儿的念头。这是生物-心理-社会因素共同作用的结果，不是矫情或不够坚强！中西医结合治疗：（1）心理治疗（CBT、人际心理治疗）；（2）必要时规范药物治疗，哺乳期用药需医师评估；（3）中药（甘麦大枣汤、逍遥散加减）需辨证使用；（4）社会支持（家庭支持、母婴关爱）。ACOG建议孕期和产后进行抑郁、焦虑等心理健康筛查；如出现自杀/伤婴念头，应立即就医或寻求紧急帮助。\end{quote}
\noindent\textit{资料依据见本页脚注。}
\footnote{佚名. 金匮要略·妇人杂病脉证并治第二十二［M］. 汉代. 电子整理本: 中国哲学书电子化计划（CTEXT）［EB/OL］. \href{https://ctext.org/jinkui-yaolue/22/zh}{在线链接}.}
\footnote{American College of Obstetricians and Gynecologists. Screening and diagnosis of mental health conditions during pregnancy and postpartum［J］. Obstetrics \& Gynecology, 2023, 141(6): 1232-1261. DOI：\nolinkurl{10.1097/AOG.0000000000005200}.}
\footnote{American College of Obstetricians and Gynecologists. Treatment and management of mental health conditions during pregnancy and postpartum［J］. Obstetrics \& Gynecology, 2023, 141(6): 1262-1288. DOI：\nolinkurl{10.1097/AOG.0000000000005202}.}
\footnote{Howard L M, Khalifeh H. Perinatal mental health: a review of progress and challenges［J］. World Psychiatry, 2020, 19(3): 313-327. DOI：\nolinkurl{10.1002/wps.20769}.}
\footnote{Stewart D E, Vigod S N. Postpartum depression［J］. New England Journal of Medicine, 2019, 380: 2187-2196. DOI：\nolinkurl{10.1056/NEJMcp1607649}.}

\subsubsection{第十二拍·汉水归舟}\label{question:zhouzhiruo:12}
\textbf{情节位置：}\hyperref[plot:zhouzhiruo:12]{汉水归舟完整剧情}。

\textbf{题干：}产后护理（坐月子）是中医文化的重要组成部分。生化汤被誉为产后第一汤。以下关于产后护理的说法，哪一项是正确的？

\textbf{选项与答案}
\begin{enumerate}[label=\Alph*.,leftmargin=*]
  \item 坐月子期间完全不能洗头和洗澡
  \item 产后可以保持清洁、尽早适度活动、均衡饮食，并按计划接受产后随访；有发热、大出血或明显情绪异常应及时就医\textbf{（正确答案）}
  \item 产后抑郁是矫情的表现，不需要在意
  \item 产后妈妈只需关注婴儿，不需要关注自身恢复
\end{enumerate}
\textbf{判定依据、推导与科普边界}
\begin{quote}现代产后照护强调连续、主动的随访，而不是把产妇关在房中忍耐。可以洗头洗澡并保持会阴和伤口清洁；无禁忌者应尽早适度活动、均衡饮食和补充水分。产后随访应在出院后尽早建立联系，并在产后数周内完成综合评估，具体时间依当地规范和个人风险安排。持续高热、胸痛气促、单侧小腿肿痛、剧烈头痛、伤口异常、突然大量出血或自伤/伤婴念头均须紧急求助。生化汤是传统方剂，若考虑使用必须先排除产后出血、感染等急症并由中医师辨证，不能作为所有产妇的常规排恶露汤。\end{quote}
\noindent\textit{资料依据见本页脚注。}
\footnote{傅山. 傅青主女科［M］. 清代刻本. 电子整理本: 识典古籍［EB/OL］. \href{https://www.shidianguji.com/book/YCM294300096/chapter/1l9ggh6dkjmer}{在线链接}.}
\footnote{World Health Organization. WHO recommendations on maternal and newborn care for a positive postnatal experience［M］. Geneva: WHO, 2022. \href{https://www.who.int/publications/i/item/9789240045989}{在线链接}.}
\footnote{American College of Obstetricians and Gynecologists. Optimizing postpartum care［J］. Obstetrics \& Gynecology, 2018, 131(5): e140-e150. DOI：\nolinkurl{10.1097/AOG.0000000000002633}.}

\subsubsection{第十三拍·峨眉金顶}\label{question:zhouzhiruo:13}
\textbf{情节位置：}\hyperref[plot:zhouzhiruo:13]{峨眉金顶完整剧情}。

\textbf{题干：}骨质疏松被称为沉默的疾病——早期无症状，却在某天一摔就骨折。以下关于骨质疏松的预防，哪一项不完全正确？

\textbf{选项与答案}
\begin{enumerate}[label=\Alph*.,leftmargin=*]
  \item 应优先从饮食获得足量钙和蛋白质；是否补充钙或维生素D需结合摄入、日照、化验及骨折风险
  \item 只有老年人才需要关注骨质疏松，年轻人无所谓\textbf{（正确答案）}
  \item 骨密度（DXA）检查应依据年龄、脆性骨折史和风险评估安排，而不是所有绝经女性每年检查
  \item 负重、抗阻和平衡训练有助于维持骨量、肌力并降低跌倒风险
\end{enumerate}
\textbf{判定依据、推导与科普边界}
\begin{quote}骨质疏松预防应贯穿全生命周期：青少年期建立较高峰值骨量，成年后保持营养、运动和健康体重，绝经后及老年期加强骨折风险评估和防跌倒。钙和维生素D的目标量、是否使用补充剂以及DXA检查时机，均应结合年龄、饮食、日照、既往骨折、用药和其他风险因素决定；不能要求所有人固定剂量补充，也不必让所有绝经女性每年做DXA。确诊高风险者可能需要规范抗骨质疏松药物。中医药如作为辅助，也应在专业指导下使用并监测安全性。\end{quote}
\noindent\textit{资料依据见本页脚注。}
\footnote{中华医学会骨质疏松和骨矿盐疾病分会. 原发性骨质疏松症诊疗指南（2022）［J］. 中国全科医学, 2023, 26(14): 1671-1691. DOI：\nolinkurl{10.12114/j.issn.1007-9572.2023.0121}.}
\footnote{国家卫生健康委员会. 中国居民骨质疏松症流行病学调查报告（2018）［R］. 北京: 国家卫生健康委员会, 2018. Available at: \href{https://www.nhc.gov.cn/ewebeditor/uploadfile/2018/10/20181020192813503.pdf}{在线链接}.}
\footnote{佚名. 黄帝内经·素问·痿论［M］. 汉代. 电子整理本: 中国哲学书电子化计划（CTEXT）［EB/OL］. \href{https://ctext.org/huangdi-neijing/wei-lun/zh}{在线链接}.}

\subsubsection{第十四拍·药王谷中}\label{question:zhouzhiruo:14}
\textbf{情节位置：}\hyperref[plot:zhouzhiruo:14]{药王谷中完整剧情}。

\textbf{题干：}四物汤（当归、白芍、熟地、川芎）被誉为妇科第一方。以下关于四物汤的说法，哪一项是正确的？

\textbf{选项与答案}
\begin{enumerate}[label=\Alph*.,leftmargin=*]
  \item 四物汤对所有女性都适用，可以在经期随意服用
  \item 四物汤的基本配伍是：当归补血活血、熟地滋阴补血、白芍养血柔肝、川芎活血行气\textbf{（正确答案）}
  \item 四物汤中的当归主要作用是清热解毒
  \item 四物汤只适用于老年女性
\end{enumerate}
\textbf{判定依据、推导与科普边界}
\begin{quote}四物汤源自唐代《仙授理伤续断秘方》，宋代《太平惠民和剂局方》将其列为妇科通用方。四味药的精妙配伍：（1）当归（君）——补血活血、调经止痛；（2）熟地（臣）——滋阴补血、益精填髓；（3）白芍（佐）——养血柔肝、缓急止痛；（4）川芎（使）——活血行气、祛风止痛。全方补血而不滞血，活血而不伤正。临床上根据具体情况化裁：血热者熟地易生地（生四物汤）；血瘀者白芍易赤芍，加桃仁红花（桃红四物汤）；气虚者加人参黄芪（圣愈汤）；兼寒者加肉桂艾叶（艾附暖宫丸）。注意事项：（1）感冒发热、脾胃湿热时不宜；（2）经期量多者慎用（川芎活血）；（3）需遵医嘱辨证使用，不可盲目自行服用。\end{quote}
\noindent\textit{资料依据见本页脚注。}
\footnote{佚名. 太平惠民和剂局方［M］. 宋代. 电子整理本: 维基文库［EB/OL］. \href{https://zh.wikisource.org/zh-hans/太平惠民和劑局方}{在线链接}.}
\footnote{卫生福利部中医药司. 四物汤说明资料［EB/OL］. \href{https://dep.mohw.gov.tw/DOCMAP/fp-866-5530-108.html}{在线链接}.}
\footnote{Lee S Y, et al. Therapeutic applications of Si-Wu-Tang in gynecological disorders: a pharmacological perspective［J］. Journal of Chinese Medicine, 2020, 31(1): 6-16.}

\subsubsection{第十六拍·终南古墓}\label{question:zhouzhiruo:16}
\textbf{情节位置：}\hyperref[plot:zhouzhiruo:16]{终南古墓完整剧情}。

\textbf{题干：}中医美容强调整体调理——以内养外。以下关于女性皮肤健康与美容的说法，哪一项是不合理的？

\textbf{选项与答案}
\begin{enumerate}[label=\Alph*.,leftmargin=*]
  \item 充足睡眠和保持心情舒畅有助于皮肤健康
  \item 追求快速美白而使用含汞、激素等违禁成分的美白产品是安全的\textbf{（正确答案）}
  \item 中医美容也需辨证与安全评估；玉容散等传统外用方并非人人适用，使用前应注意刺激和过敏风险
  \item 防晒是延缓皮肤光老化的最重要外用措施
\end{enumerate}
\textbf{判定依据、推导与科普边界}
\begin{quote}非法添加汞化合物、糖皮质激素等的速效美白产品可能造成肾脏、神经系统或皮肤屏障损害；来源不明产品还可能存在重金属超标。更可靠的皮肤健康策略是防晒、温和清洁、按肤质保湿和治疗明确的皮肤病。防晒应结合遮阳、衣帽与足量广谱防晒产品。传统外用方和草本成分也可能刺激或致敏；四物汤、逍遥散、六味地黄丸等内服方不是美容保健品，不应自行长期服用。\end{quote}
\noindent\textit{资料依据见本页脚注。}
\footnote{国家药品监督管理局. 化妆品安全技术规范（2015年版）及相关修订通告［S/OL］. \href{https://www.nmpa.gov.cn/}{在线链接}.}
\footnote{World Health Organization. Radiation: protecting against skin cancer［EB/OL］. \href{https://www.who.int/news-room/questions-and-answers/item/radiation-protecting-against-skin-cancer}{在线链接}.}
\footnote{Krutmann J, Bouloc A, Sore G, et al. The skin aging exposome［J］. Journal of Dermatological Science, 2017, 85(3): 152-161. DOI：\nolinkurl{10.1016/j.jdermsci.2017.09.015}.}
\footnote{吴谦. 医宗金鉴·外科卷上［M］. 清代刻本. 电子整理本: 维基文库［EB/OL］. \href{https://zh.wikisource.org/zh-hans/醫宗金鑑/外科卷上}{在线链接}.}
\footnote{中国中医药出版社. 中医美容学（第2版）［M］. 北京: 中国中医药出版社, 2021.}

\subsubsection{第十七拍·华山论剑}\label{question:zhouzhiruo:17}
\textbf{情节位置：}\hyperref[plot:zhouzhiruo:17]{华山论剑完整剧情}。

\textbf{题干：}宫颈癌是目前最有希望通过HPV疫苗、规范筛查和早治走向公共卫生消除的癌症之一（WHO 2020年提出全球消除宫颈癌战略）。以下关于HPV疫苗和宫颈癌筛查的说法，哪一项是不正确的？

\textbf{选项与答案}
\begin{enumerate}[label=\Alph*.,leftmargin=*]
  \item HPV疫苗在首次性行为前接种获益最大；超过重点接种年龄后是否接种，应依据当地免疫政策、年龄和个人情况咨询专业人员
  \item 接种HPV疫苗后不需要再进行宫颈癌筛查\textbf{（正确答案）}
  \item 宫颈癌筛查的起始年龄、检测方法和间隔应按所在地指南、既往结果与个人风险决定
  \item 宫颈癌是可以预防和治愈的——前提是早筛查、早发现、早治疗
\end{enumerate}
\textbf{判定依据、推导与科普边界}
\begin{quote}接种HPV疫苗后仍需要按指南参加宫颈癌筛查。现有疫苗可预防多数但并非全部致癌型别，也不能治疗接种前已经存在的感染；筛查则用于发现高危HPV持续感染和癌前病变。具体起始年龄、采用HPV检测还是细胞学/联合检测、间隔多久，以及何时停止筛查，均受所在地指南、免疫状态、既往结果和个人风险影响，不能统一写成所有人每3-5年一次。WHO消除宫颈癌战略以疫苗、筛查和规范治疗三条路径共同推进，三者不能互相替代。\end{quote}
\noindent\textit{资料依据见本页脚注。}
\footnote{World Health Organization. Cervical cancer fact sheet[EB/OL]. 2025. \href{https://www.who.int/news-room/fact-sheets/detail/cervical-cancer}{在线链接}.}
\footnote{World Health Organization. Global strategy to accelerate the elimination of cervical cancer as a public health problem[S]. 2020. \href{https://www.who.int/publications/i/item/9789240014107}{在线链接}.}
\footnote{Centers for Disease Control and Prevention. HPV vaccine recommendations[EB/OL]. 2024. \href{https://www.cdc.gov/hpv/hcp/vaccination-considerations/index.html}{在线链接}.}
\footnote{Lei J, Ploner A, Elfström K M, et al. HPV vaccination and the risk of invasive cervical cancer[J]. New England Journal of Medicine, 2020, 383(14):1340–1348. DOI：\nolinkurl{10.1056/NEJMoa1917338}.}
\footnote{China CDC Weekly. Accelerating cervical cancer prevention and control in China to achieve elimination strategy objectives[J]. 2022, 4(48):1067–1069. DOI：\nolinkurl{10.46234/ccdcw2022.215}.}

\subsubsection{第十八拍·白莲重生}\label{question:zhouzhiruo:18}
\textbf{情节位置：}\hyperref[plot:zhouzhiruo:18]{白莲重生完整剧情}。

\textbf{题干：}纵观周芷若十八关的历程——从营养不良的渔家女到独当一面的峨眉掌门。以下哪一项最能总结女性健康的核心理念？

\textbf{选项与答案}
\begin{enumerate}[label=\Alph*.,leftmargin=*]
  \item 女性健康只是妇科问题
  \item 治未病理念——预防优于治疗，女性全生命周期健康管理需要生理、心理、社会三位一体的综合关怀\textbf{（正确答案）}
  \item 女性只要身体健康就好，不需要关注心理健康
  \item 女性健康问题随着年龄增长自然会解决
\end{enumerate}
\textbf{判定依据、推导与科普边界}
\begin{quote}周芷若的十八关涵盖了女性健康的多个维度：青春期发育（第二拍）、月经健康（第二/四拍）、精神压力与月经（第四拍）、药物安全（第五拍）、情感创伤（第六拍）、运动损伤（第七拍）、职业女性健康（第八拍）、失眠与焦虑（第九拍）、自我价值（第十拍）、产后护理（第十二拍）、更年期与骨质疏松（第十三拍）、乳腺健康（第十五拍）、宫颈癌预防（第十七拍）——这正是女性全生命周期健康管理的完整图景。中医治未病思想——圣人不治已病治未病，不治已乱治未乱——与WHO倡导的预防医学理念高度契合。女性健康≠妇科，而是涵盖心血管健康、骨骼健康、精神心理健康、肿瘤防治、性生殖健康等全方位的系统工程。愿每一位女性——如周芷若最终领悟的那样——先爱自己，再爱他人。\end{quote}
\noindent\textit{资料依据见本页脚注。}
\footnote{中国哲学书电子化计划. 黄帝内经·素问·四气调神大论[EB/OL]. \href{https://ctext.org/huangdi-neijing/si-qi-diao-shen-da-lun/zhs}{在线链接}.}
\footnote{World Health Organization. Women’s health[EB/OL]. \href{https://www.who.int/health-topics/women-s-health}{在线链接}.}
\footnote{Pan American Health Organization; World Health Organization. Healthy life course[EB/OL]. \href{https://www.paho.org/en/topics/healthy-life-course}{在线链接}.}
\footnote{中华人民共和国国务院. 健康中国2030规划纲要[S]. 2016. \href{https://www.gov.cn/zhengce/2016-10/25/content_5124174.htm}{在线链接}.}
\footnote{中华人民共和国国务院. 中国妇女发展纲要（2021—2030年）[S]. 2021.}
